\section{Introduction}

Forest management in British Columbia is decided at two levels that do not share an objective.
The Chief Forester sets an allowable annual cut for each timber supply area, with the duty to
sustain the timber supply over time~\citep{government_of_british_columbia_aac_2026}. The
enterprises that hold the harvesting rights then decide where, when, and with which species to
harvest and regenerate, answering to price, cost, and the terms of their tenure. Policy wants a
forest that stays productive and resilient under fire risk. The enterprise wants a margin. When
the two objectives diverge, a policy that ignores the enterprise response misses its target, and
an enterprise that ignores policy is corrected after the fact, at a cost.

The divergence matters more as the disturbance regime changes. Forest disturbances in Europe have
increased in extent and severity over recent decades~\citep{seidl_increasing_2014}, and the
Canadian forest sector faces the same direction of change, with fire seasons that exceed the
historical range and a timber supply that carries the legacy of beetle-killed pine and past
harvest. Under such a regime, the composition of the forest is not only a silvicultural choice but
a risk decision: species differ in how they burn, how they grow, and what they are worth.

Tree-species diversification has been proposed as the adaptation strategy that reconciles the two
objectives. The economic case was tested at scale for European forests by
\citet{fuchs_severe_2024}, who optimized stand-type portfolios across planning units under both
small-scale disturbances and spatially correlated extreme weather events. Diversification raised
expected returns and reduced risk from isolated disturbances, but severe and frequent extreme
events undermined the economic gain, and the optimal response shifted toward homogenisation and
low-investment species. That result is the point of departure for this paper: the benefit of
mixing is conditional on the disturbance regime, and a risk analysis that treats disturbances as
independent draws will overstate it.

Two further lines of work bound the question. \citet{topanotti_enhancing_2024} showed empirically
that enrichment of production forests can serve economic and ecological multifunctionality at the
same time, so composition is not a simple trade-off between profit and diversity. And
\citet{zelnik_managing_2025} showed with a process model that managing for heterogeneity, in stand
size and in conifer--deciduous composition, reduces fire risk at the landscape scale. Together they
suggest that the relevant composition decision is spatial and dynamic, not a single landscape
average.

What is missing is a study that joins these pieces. The European portfolio analysis represents
disturbance as a probability over stylized units and allocates area shares of stand types; it does
not simulate fire spread, and its units do not have the age structure, access, and adjacency of an
operating forest. The fire-heterogeneity work simulates spread but has no economic agent, no
harvest scheduling, and no policy constraint. Neither asks what an enterprise would actually do
under a policy that targets composition, nor what that response costs in the tail of the
distribution rather than at the mean.

This paper builds that bridge for one real management unit, the Prince George Timber Supply Area
(TSA~24) in interior British Columbia. We formulate a bilevel problem in which a regional policy
maker chooses an allowable annual cut and a target species composition, while an enterprise
responds with a harvest and regeneration schedule that maximizes net present value subject to
those choices. The estate is the compiled TSA~24 inventory, with its actual blocks, ages, species,
and yield curves. Fire is a spatially explicit process coupled to the Canadian Forest Fire
Behaviour Prediction system, calibrated against the historical perimeter record for the area, and
driven by reconstructed weather; it enters the analysis as a scenario, and policies are compared
against each other on identical fire draws. Risk is measured as the conditional value at risk of
net present value, following \citet{fuchs_severe_2024}, so that the tail of the distribution, not
only its mean, is part of the objective.

We ask four questions:

\begin{enumerate}
    \item[\textit{RQ1}] Which landscape-scale species mixtures best balance economic production
    capacity and fire risk under a stated fire regime?
    \item[\textit{RQ2}] What are the synergies and trade-offs between species diversification and
    the rotation length implied by the allowable cut?
    \item[\textit{RQ3}] How do the economically favourable mixtures change with the decision
    maker's risk attitude?
    \item[\textit{RQ4}] Over which planning horizons and discount rates does mixing become
    economically favourable?
\end{enumerate}

A fifth question motivates the work but is not answered here: to what extent non-productive
broadleaf species can be placed as fire breaks before production is over-proportionally
compromised. Answering it requires a spatial placement rule for the breaks, which the policy level
in this model does not have; we return to it as a limitation and as the natural extension of the
framework.

The contribution is therefore threefold. First, a bilevel formulation that makes the policy and
enterprise objectives explicit on a real estate, rather than assuming that area shares translate
into outcomes. Second, a coupling of that formulation to a calibrated, spatially explicit fire
process, so that fuel structure, composition, and connectivity affect the risk that the policy is
meant to manage. Third, an evaluation on downside risk, reported with the sampling uncertainty of
the historical fire record, which is what makes the comparison between policies meaningful when
the tail of the disturbance distribution is itself weakly identified.

The paper proceeds as follows. \autoref{sec:methods} describes the study area, the two decision
levels, the estate model, the fire scenario, and the economic valuation.
\autoref{sec:results} reports the estate validation, the even-flow schedule against the official
cut, the economic results and their sensitivities, and the fire comparison.
\autoref{sec:discussion} interprets the results against the European evidence on diversification,
states the limitations, and closes with what a policy maker should take from the analysis.
